\documentclass[10pt,letterpaper]{article}
\usepackage[top=0.85in,left=2.75in,footskip=0.75in]{geometry}
\usepackage{amsmath,amssymb}
\usepackage{changepage}
\usepackage[utf8x]{inputenc}
\usepackage{textcomp,marvosym}
\usepackage{cite}
\usepackage{nameref,hyperref}
\usepackage[right]{lineno}
\usepackage{microtype}
\DisableLigatures[f]{encoding = *, family = * }
\usepackage[table]{xcolor}
\usepackage{array}
\usepackage{listings}
\usepackage{graphicx}
\usepackage{float}
\usepackage{booktabs}
\usepackage{multirow}

% Note: This template is based on PLOS Computational Biology requirements.
% It assumes you have the plos2015.bst file for bibliography if you use \bibliographystyle{plos2015}.
% If not, use \bibliographystyle{unsrt} or similar.

\linenumbers

\title{MedImages.jl: A High-Performance, Differentiable Julia Framework for Medical Image Processing and Metadata Preservation}

\author{Jakub Mitura$^{1*}$, Divyansh Goyal$^{1}$, Joanna Wybranska$^{1}$, Michael C. Kreissl$^{1}$}

\date{}

\begin{document}

\maketitle

\noindent
$^{1}$ Division of Nuclear Medicine, Department of Radiology \& Nuclear Medicine, Faculty of Medicine, Otto von Guericke University Magdeburg, 39120 Magdeburg, Germany \\
$^{*}$ Corresponding author: jakub.mitura@med.ovgu.de \\
Full Postal Address: Leipziger Str. 44, 39120 Magdeburg, Germany \\
Telephone: [Insert Telephone Number Here]

\section*{Abstract}
\textbf{Background and Objectives:} Quantitative medical imaging, particularly in Nuclear Medicine, requires software that ensures both high-throughput computational performance and strict spatial metadata integrity. Existing frameworks often face a "two-language" bottleneck, bridging high-level APIs with compiled backends, which complicates automatic differentiation (AD) and can lead to metadata drift during preprocessing. We introduce MedImages.jl, a native Julia ecosystem for 3D/4D medical image analysis.

\textbf{Methods:} MedImages.jl implements a typed data structure that intrinsically binds voxel data to spatial metadata (origin, spacing, direction). Core geometric transformations are implemented as native, differentiable Julia kernels compatible with GPU acceleration (KernelAbstractions.jl) and SciML AD engines (Zygote.jl). The framework integrates HDF5 for high-throughput I/O and supports batched multimodality operations.

\textbf{Results:} In a 100-subject PET/CT preprocessing benchmark, MedImages.jl demonstrated a 7.2$\times$ total turnaround speedup compared to standard Python-based caching pipelines. GPU-accelerated spatial transformations achieved up to 135$\times$ speedups over established single-threaded C++ baselines like SimpleITK. We validated end-to-end differentiability through a spatial transformer network and multi-organ affine registration, and confirmed that Standardized Uptake Value (SUV) calculations match established clinical tools to double-precision accuracy.

\textbf{Conclusions:} MedImages.jl provides a performant, fully differentiable environment for medical image processing. By unifying spatial operations and metadata management within a single language, it facilitates large-scale clinical research and physics-informed deep learning without compromising clinical validity.

\textbf{Keywords:} Medical Imaging, Julia, Scientific Machine Learning, Nuclear Medicine, Differentiable Programming, Spatial Metadata

\section{Introduction}

The development of MedImages.jl is driven by the need to address four converging challenges in contemporary medical imaging pipelines, particularly within Nuclear Medicine workflows that rely heavily on functional data (e.g., PET, SPECT) and anatomical references (CT, MRI).

\paragraph{Challenge 1: The Volume of Modern Biobank Datasets.}
A single whole-body PET/CT examination involves high-resolution 3D volumes that must be co-registered, resampled, and quantitatively normalized. At biobank scale---encompassing thousands of subjects---preprocessing these massive volumetric datasets becomes a dominant bottleneck. In traditional hybrid environments (SimpleITK \cite{LowekampChen2013}, MONAI \cite{CardosoLi2022}), I/O overhead and metadata synchronization account for a significant portion of processing time, leading to substantial delays in multimodal dataset preparation \cite{Dexl_2025}.

\paragraph{Challenge 2: Execution Speed and the "Two-Language" Barrier.}
Researchers often prototype in high-level languages like Python but must rely on compiled C++ backends for execution speed. This "two-language problem" complicates deployment and limits code inspectability. Julia \cite{BezansonChen2018} addresses this by compiling directly to efficient machine code via LLVM JIT, allowing spatial transformations to be executed rapidly across both CPU and GPU backends using native implementations rather than wrapped binaries.

\paragraph{Challenge 3: Differentiability and the Fragmented SciML Landscape.}
Emerging Scientific Machine Learning (SciML) paradigms require pipelines where every computational step is differentiable. While Python offers powerful differentiable libraries (e.g., PyTorch, JAX, TensorFlow), they often function as "walled gardens." Composing a pipeline requires forcing the entire ecosystem to commit to a single framework, creating friction when integrating third-party tools or classical numerical solvers \cite{PalBhattacharya2024, EschleGl2023}. Julia addresses this structural fragmentation by solving the "expression problem" through multiple dispatch, allowing native code across different packages to seamlessly interoperate \cite{RoeschGreener2023}. Because Julia's automatic differentiation engines (e.g., Zygote.jl \cite{Fischer_2019}, Enzyme.jl) operate at the compiler level, they can differentiate through arbitrary code, including complex differential equation solvers and GPU kernels, without requiring users to rewrite logic into a specific tensor framework. MedImages.jl introduces a critical tool to make this unique, highly composable SciML ecosystem widely accessible to medical imaging researchers.

\paragraph{Challenge 4: Metadata Management and Physical Interpretation.}
Unlike standard computer vision, medical images possess a strict "physical interpretation." A medical image is a grid of points occupying a physical region defined by an origin, spacing, and direction. The transition from medical image formats (DICOM, NIfTI) to raw numerical arrays for deep learning often strips this spatial context. This risks "metadata drift"—a potentially dangerous misalignment between voxel data and its spatial coordinates during preprocessing or augmentation.

MedImages.jl provides a unified, GPU-accelerated, and differentiable imaging framework designed explicitly to address these four challenges. The structure of this paper, including the methodology and results, is organized around validating solutions to these core axes.

\subsection{Architectural Foundations: The Coordinate System and Metadata}
Addressing Challenge 4 (Metadata Management) requires ensuring that data remains defined within a physical coordinate system. To accurately integrate functional data with anatomical data across different spatial resolutions, software must rigorously maintain the mapping between discrete voxel indices and physical patient space. This mapping from voxel coordinates $\mathbf{v} = [i, j, k, 1]^\top$ to continuous world coordinates $\mathbf{p} = [x, y, z, 1]^\top$ relies on a $4 \times 4$ homogeneous affine transformation matrix $M$:
\begin{equation}
\mathbf{p} = M \mathbf{v}, \quad M = 
\begin{bmatrix} 
R S & T \\
\mathbf{0} & 1
\end{bmatrix}
\end{equation}
where $R$ is the $3 \times 3$ rotation matrix, $S$ is a diagonal scaling matrix (voxel spacing), and $T$ is the translation vector (origin). MedImages.jl ensures that any spatial operation $\mathcal{T}$ applied to the voxel grid is mathematically coupled with an update $M_{new} = M_{old} \times M_{\mathcal{T}}$, preventing spatial desynchronization.

Furthermore, quantitative nuclear medicine relies on precise temporal and clinical metadata. The Standardized Uptake Value (SUV), a critical metric for PET/CT, is calculated as:
\begin{equation}
SUV = \frac{C(t)}{ID / BW}
\end{equation}
where $C(t)$ is the decay-corrected radioactive concentration, $ID$ is the injected dose, and $BW$ is patient body weight. Discarding parameters such as \textit{Acquisition Time} or \textit{Patient Weight} compromises the quantitative utility of the image. MedImages.jl intrinsically protects these fields within its data structures.

\section{Methods}

\subsection*{Software Architecture}
MedImages.jl implements a metadata paradigm inspired by the Brain Imaging Data Structure (BIDS) standard, organizing and validating clinical metadata rigorously. The core \texttt{MedImage} struct enforces metadata integrity through Julia's type system. Unlike array-based deep learning workflows where converting to tensors discards spatial context, MedImages.jl intrinsically binds metadata (origin, spacing, direction, patient ID, acquisition timestamp) to the structural definition of the medical image.

ITK (C++) is used via wrapper primarily for I/O robustness (NIfTI, DICOM parsing), while spatial transformations are implemented natively in Julia. This ensures: (1) \textbf{Differentiability}—operations integrate with Zygote.jl/Enzyme.jl; (2) \textbf{Composability}—seamless integration with Lux.jl and DifferentialEquations.jl.

To support high-throughput analysis of multimodal studies, MedImages.jl introduces \texttt{BatchedMedImage}, which stores 4D voxel data $(x, y, z, \text{batch})$. Crucially, the spatial metadata (origin, spacing, direction) is vectorized. This design allows each 3D slice within the batch to maintain unique spatial properties while residing in a unified 4D tensor, facilitating efficient batched operations where identical geometric transformations can be dispatched to the GPU in a single pass. To minimize I/O bottlenecks for large datasets, the framework supports HDF5-based storage.

\subsubsection*{Code Example: Typical Workflow}
The following snippet demonstrates loading a multimodal pair, rotating them simultaneously using the batched structure, and computing the SUV:

\begin{lstlisting}[basicstyle=\ttfamily\footnotesize]
using MedImages

# Load PET and CT (DICOM metadata is preserved)
pet = load_image("patient_1/pet.dcm")
ct = load_image("patient_1/ct.nii")

# Create a batched tensor to process simultaneously
batch = create_batched_medimage([pet, ct])

# Rotate both volumes by 45 degrees around Z-axis (GPU supported)
rotated_batch = rotate_mi(batch, 45.0)

# Extract SUV statistics for the PET scan using an ROI mask
suv_stats = calculate_suv_statistics(rotated_batch[1], roi_mask)
println("Mean SUV: ", suv_stats.mean_suv)
\end{lstlisting}

\subsection*{Functionalities and I/O Robustness}
MedImages.jl provides a comprehensive set of functionalities optimized for 3D medical image processing:
\begin{itemize}
    \item \textbf{Robust I/O:} Loading of various medical image formats (NIfTI, DICOM) using \texttt{ITKIOWrapper}.
    \item \textbf{Native Julia Transformations:} Functions to modify spatial properties of 3D volumes (\texttt{rotate\_mi}, \texttt{translate\_mi}, \texttt{scale\_mi}, \texttt{resample\_to\_spacing}, \texttt{crop\_mi}, \texttt{pad\_mi}).
    \item \textbf{Resample to Target:} Resamples a moving image to match the spatial geometry of a reference image.
    \item \textbf{GPU Acceleration:} All transformations transparently support GPU execution via KernelAbstractions.jl.
    \item \textbf{HDF5 Support:} Fast loading and saving of 3D images using HDF5, speeding up I/O for intermediate processing and reducing the "loading penalty" common in deep learning pipelines.
    \item \textbf{Autodifferentiation:} All spatial transformations have well-defined gradients via Zygote.jl and Enzyme.jl.
    \item \textbf{Orientation Management:} Robust tools to handle and convert between neuroimaging coordinate systems.
\end{itemize}

\subsection*{Experimental Setup: Evaluating the Four Challenges}
To systematically evaluate the framework against the four core challenges, several experimental pipelines were designed.

\subsubsection*{Evaluating Volume and Speed}
For large-scale studies requiring standardized preprocessing of heterogeneous medical volumes, we developed a fused affine preprocessing pipeline that collapses multiple spatial transformations---orientation normalization, spacing resampling, and spatial centering---into a single $4 \times 4$ affine matrix. This matrix is precomputed on the CPU from spatial metadata, then applied via a custom KernelAbstractions.jl kernel with unrolled trilinear interpolation on the GPU.

To evaluate the framework under high-throughput production loads (Challenge 1: Volume), we benchmarked a 100-case preprocessing pipeline on the autoPET dataset \cite{Dexl_2025} against MONAI's \texttt{PersistentDataset} mechanism. A 100-case dataset was selected to represent a typical multimodality workflow bottleneck: physically disparate modalities (PET and CT) must be loaded and strictly aligned to a common spatial reference frame---meaning they share the identical voxel spacing, direction cosines, and origin---so they can be safely concatenated into a single multi-channel tensor for deep learning ingestion. 

To assess raw execution speed (Challenge 2: Speed), we benchmarked isolated transformations across multiple scales. SimpleITK benchmarks were executed strictly on the CPU, as it does not natively support GPU acceleration for these geometric filters. Conversely, MedImages.jl and MONAI benchmarks utilized GPU execution to demonstrate the high throughput enabled by natively fused kernels. All experiments described are open-source and reproducible via the codebase provided in this repository.

\subsubsection*{Evaluating Differentiability (SciML Integrations)}
To validate end-to-end differentiability (Challenge 3), we implemented three paradigms:
\begin{enumerate}
    \item \textbf{Learned Inverse Rotation (Differentiable Geometric CNN):} A neural network learns to predict and invert unknown 3D rotations applied to volumetric data. To avoid confusion with modern attention-based architectures, we clarify that this uses a standard 3D Convolutional Neural Network (CNN) as a spatial regressor. The CNN predicts three Euler angles, which are converted to an affine matrix to generate a differentiable sampling grid for MedImages.jl's \texttt{interpolate\_pure} function. 
    \item \textbf{Organ-Specific Affine Registration:} A 3D CNN predicts per-organ affine parameters. The loss computation is a single GPU kernel fusing affine transformation, a barycenter distance penalty, and trilinear interpolation loss, differentiated via \texttt{Enzyme.autodiff}.
    \item \textbf{Quantitative Dosimetry Refinement via Universal Differential Equations:} A comprehensive framework using \texttt{DifferentialEquations.jl} \cite{Rackauckas_2017} to formulate a 4-state Universal Differential Equation (UDE) for high-fidelity $^{177}$Lu-PSMA dosimetry. This approach corrects approximate analytical clinical models (like Voxel S-Values, VSV) into Monte Carlo-level patient-specific dose maps by embedding a neural network corrector $\mathcal{N}_\theta$ directly within the radiation transport and pharmacokinetic kinetic systems.
\end{enumerate}

The clinical imperative for high-accuracy dosimetry has never been more acute as radiation oncology transitions toward personalized Targeted Radionuclide Therapy (TRT), such as $^{177}$Lu-PSMA. Traditional clinical workflows rely on simplified analytical methods---specifically, Voxel S-Value (VSV) convolution over Time-Integrated Activity (TIA) maps. This methodology uses mono-exponential fitting across longitudinal SPECT time-points and assumes a homogeneous water-equivalent medium ($S_{homo}$ kernels), leading to significant errors at tissue interfaces or highly heterogeneous density regions (e.g., lungs, bone cortex). While Monte Carlo (MC) techniques solve this by tracking complex stochastic particle interactions (Compton/Rayleigh scattering, Bremsstrahlung), their immense computational burden heavily restricts their clinical utility. Conversely, purely deep learning-based approaches often lack the necessary physical constraints and safety guarantees.

To resolve this dichotomy, our pipeline models dose deposition and pharmacokinetic (PK) binding as a parameterized dynamical system. By explicitly partitioning the dosimetry problem into \textit{Mechanistic Knowns} and \textit{Learned Uncertainties}, UDEs satisfy the critical need for interpretability. The system treats the absorbed dose rate as a discoverable residual within a robust physical framework, solved as an Initial Value Problem (IVP).

\paragraph{Mechanistic Constraints (The Knowns):}
The UDE model heavily constrains the solution space using established physics and biological priors \cite{KennedyLugassi2020}. The compartmental activity kinetics depend on the physical decay constant $\lambda_{phys} = 4.34 \times 10^{-3} \text{ h}^{-1}$ \cite{ReschFouladgar2023, StabinMadsen2019}, the blood circulation rate ($\lambda_{bc} \approx 0.693 \text{ min}^{-1}$), and the renal clearance scaling with the patient's Glomerular Filtration Rate (GFR) \cite{HardiansyahYousefzadehNowshahr2024}.

Crucially, the mechanistic model rigorously integrates quantitative imaging transformations. The observed SPECT activity $A_{obs}(\mathbf{r}, 0)$ is corrected via the Camera Calibration Factor ($CF$) derived from DICOM tags \cite{BidgoodHorii1992, GleisnerChouin2022}, and spatial Recovery Coefficients ($RC$) \cite{KlettingSchimmel2015} before integration:
\begin{equation}
    A_{true}(\mathbf{r}, 0) = \frac{A_{obs}(\mathbf{r}, 0) \cdot RC(\mathbf{r})}{CF}
\end{equation}
For radiation transport, the baseline dose rate $\dot{D}_{base}$ is computed via 3D FFT convolution of the total activity against $S_{homo}$ kernels \cite{PistoneAuditore2022}. To handle heterogeneous tissues correctly, the target voxel mass $m(\mathbf{r})$ is derived from the anatomical CT. A bi-linear calibration curve converts Hounsfield Units (HU) to a physical mass density map $\rho(\mathbf{r})$ \cite{GtzSchmidkonz2019}, allowing proper energy-to-dose normalization ($D \propto \text{Energy} / m(\mathbf{r})$) \cite{StabinMadsen2019, BaileyMarquis2022}.

\paragraph{Data-Driven Discovery (The Learned Uncertainties):}
With the primary physical and pharmacokinetic laws fixed, the universal component embedded within the system ($\mathcal{N}_\theta$) targets the highly complex phenomena that are difficult or impossible to measure directly. For the PK sub-system, this includes the local receptor saturation thresholds (Voxel-Wise $B_{MAX}$, or the Tumour Sink Effect) \cite{BegumThieme2018} and the cycle-dependent kinetic drift of uptake and washout rates ($k_{in}, k_{out}$) as tumours regress \cite{OkamotoThieme2016, SiebingaVeen2024, SiebingaPriv2023}. For the transport sub-system, $\mathcal{N}_\theta$ learns the residuals caused by heterogeneous mass-energy absorption ($\mu_{en}/\rho$) and complex scattering at density gradients ($\nabla \rho$) \cite{MikellKappadath2016, BaileyMarquis2022}.

The resulting coupled 4-state system describes the temporal evolution of blood activity ($A_{blood}$), free tissue activity ($A_{free}$), bound tissue activity ($A_{bound}$), and the absorbed dose rate $\dot{D}$:
\begin{align}
    \frac{dA_{blood}}{dt} &= -\left(k_{10} \cdot \frac{GFR}{120} + \lambda_{phys} \right) A_{blood}(t) - \sum \text{Uptake} + \sum \text{Washout} \\
    \frac{dA_{bound}}{dt} &= k_3 A_{free} \left( 1 - \frac{A_{bound}}{B_{MAX}(\mathbf{r})} \right) - (k_4 + \lambda_{phys}) A_{bound} \\
    \dot{D}(\mathbf{r}, t) &= \text{Softplus} \left( \frac{1}{m(\mathbf{r})} \left( \left[ A_{total}(t) \cdot \bar{E} \ast S_{homo} \right](\mathbf{r}) + \mathcal{N}_\theta \Big( A_{total}, \rho, \nabla \rho \Big) \right) \right)
\end{align}

To ensure strict differentiability across the entire 5D tensor domain without encountering Zygote.jl array mutability tracking limitations, this SciML IVP is solved using a custom, out-of-place manual Euler integration loop rather than standard \texttt{SciMLSensitivity} adjoints. We evaluated this "No-Approx UDE" against standard analytical baselines (VSV Evaluate-Proj) and multiple data-driven baselines (Stabilized CNN, PBPK-ML Extra Trees, Spect0 3D Residual, and DblurDoseNet).


\section{Results}
To demonstrate that MedImages.jl resolves the core computational and metadata challenges, our results are structured along the axes defined in the Introduction.

\subsection*{Challenge 1: Volume -- Scaling to 100 Cases (HDF5 vs. Caching)}
To evaluate the framework under high-throughput loads, we benchmarked a 100-case preprocessing pipeline on the autoPET dataset. The core computational task was a multi-modal alignment pipeline, resampling highly heterogeneous PET and CT volumes to a standardized $512 \times 512 \times 512$ isotropic grid. 

In clinical practice, "alignment" implies mathematically mapping the voxel grid of a functional image (e.g., PET, which often has lower resolution and varying spatial origins) to the exact coordinate space of an anatomical reference image (e.g., CT). This requires generating a continuous interpolation grid based on the source metadata (origin, spacing, and direction cosine matrix) and resampling the raw voxel intensities. This alignment is the foundational prerequisite for modern multimodality analysis; Convolutional Neural Networks (CNNs) require multi-channel tensors where pixel $(i,j,k)$ in the PET channel corresponds to the exact same physical tissue location as pixel $(i,j,k)$ in the CT channel. Due to the massive scale of 3D modalities, this continuous spatial resampling is often the dominant bottleneck in deep learning pipelines.

For a fair comparison of computational overhead in this high-volume scenario, we evaluated performance for high-frequency training loops (Epochs 2+), assuming data has been initially processed and cached. We compared MedImages.jl utilizing HDF5 against MONAI's \texttt{PersistentDataset} mechanism. It should be noted that this compares different data persistence strategies—HDF5 vs. Pickle/Pt-based caching—as much as the underlying framework performance.

\begin{table}[H]
\centering
\caption{Definitive Performance Comparison (100 Subjects, Epochs 2+)}
\begin{tabular}{|l|c|c|c|}
\hline
\textbf{Phase (per Subject)} & \textbf{MedImages (HDF5)} & \textbf{MONAI (PersistentDataset)} & \textbf{Speedup} \\
\hline
Stage 1: Load from Cache & \textbf{~40 ms} & ~150 ms & 3.7$\times$ \\
Stage 2: Transform (Compute) & \textbf{~50 ms} & ~500 ms & 10.0$\times$ \\
\hline
\textbf{Total Turnaround} & \textbf{~90 ms} & \textbf{~650 ms} & \textbf{7.2$\times$} \\
\hline
\end{tabular}
\label{tab:scaling_comparison}
\end{table}

Overall, MedImages.jl achieved a 7.2$\times$ total turnaround speedup compared to the MONAI PersistentDataset baseline. This operational throughput advantage is visually depicted in Figure \ref{fig:speed_comparison}.

\begin{figure}[H]
    \centering
    \includegraphics[width=0.7\textwidth]{elsarticle/figures_new/speed_comparison.png}
    \caption{\textbf{High-Volume Pipeline Turnaround.}
    Turnaround times per subject during high-frequency multimodality training loops (Epochs 2+). MedImages.jl via HDF5 outpaces the MONAI PersistentDataset caching mechanism, heavily reducing both the I/O cache load and transform compute penalties.}
    \label{fig:speed_comparison}
\end{figure}

\subsection*{Challenge 2: Speed -- GPU Acceleration Benchmarks}
MedImages.jl was benchmarked on CPU and GPU backends (NVIDIA RTX 3090, Intel Core Ultra 9 285K). For $256 \times 256 \times 128$ volumes, MedImages.jl on GPU achieved significant speedups compared to SimpleITK on CPU. Figure \ref{fig:transform_benchmarks} visualizes these execution time differences across core spatial operations:
\begin{itemize}
    \item \textbf{Resampling (Nearest Neighbor):} 115$\times$ speedup vs CPU, 17$\times$ vs SimpleITK.
    \item \textbf{Orientation Changes:} 71$\times$ speedup vs CPU, 31$\times$ vs SimpleITK.
    \item \textbf{Fused Affine Transformation:} 135$\times$ speedup vs CPU, 8.1$\times$ vs SimpleITK.
\end{itemize}

\begin{figure}[H]
    \centering
    \includegraphics[width=0.8\textwidth]{elsarticle/figures_new/transform_benchmarks.png}
    \caption{\textbf{Execution Time Comparison for Core Spatial Operations.}
    Log-scale plot comparing execution times (in milliseconds) for a $256 \times 256 \times 128$ volume across SimpleITK (CPU), MedImages CPU, and MedImages GPU. Native Julia GPU acceleration via \texttt{KernelAbstractions.jl} drastically lowers computational overhead, culminating in up to 135$\times$ speedups for fused affine transforms.}
    \label{fig:transform_benchmarks}
\end{figure}

Regarding the Fused Affine Pipeline defined in the methods, benchmarking on 100 subjects resampled to a $512^3$ target grid demonstrated per-subject GPU kernel times of under 200\,ms, confirming the pipeline's suitability for biobank-scale processing.

\subsection*{Challenge 3: Differentiability and SciML Validation}
The Julia ecosystem is uniquely positioned for SciML because its automatic differentiation frameworks (Zygote.jl, Enzyme.jl) can differentiate through arbitrary native code, including custom spatial transformations. Nuclear medicine imaging is inherently quantitative, but its key readouts (e.g., SUV, time–activity curves, absorbed dose) are highly sensitive to geometric preprocessing such as resampling, rotation, and registration. In practice, these operations must often be integrated into AI pipelines for motion correction, longitudinal alignment, and learned preprocessing, where transform parameters are optimized jointly with network weights. However, many commonly used imaging backends execute spatial resampling outside the automatic differentiation graph, preventing end-to-end training through geometric steps \cite{Berman_2024}. To demonstrate that MedImages.jl enables fully differentiable volumetric spatial operations suitable for nuclear medicine workflows, we implement several SciML validation tests.

\subsubsection*{Differentiability Verification Tests}
To rigorously validate correctness, we verified that Zygote.jl returns non-trivial, structurally correct gradients across all spatial transforms. For identity-preserving operations (\texttt{translate\_mi}, \texttt{pad\_mi}), the gradient propagates as $\mathbf{1}$ without corruption. For spatially selective operations (\texttt{crop\_mi}), the gradient is $1$ inside the cropped region and $0$ outside. For the fused affine interpolation kernel, we validated CPU/GPU gradient consistency to within $10^{-4}$.

\subsubsection*{Learned Spatial Transformations}
In the Learned Inverse Rotation experiment, end-to-end differentiability of the geometric pipeline was confirmed using synthetic single-case batched configurations. Over 20 training iterations, the CNN successfully optimized the predicted Euler angles, consistently reducing the mean squared reconstruction error (MSE) between the de-rotated output and the original target by over 65\%. This demonstrates that differentiable augmentation and dynamic learned preprocessing---which are often difficult or rare in standardized imaging pipelines \cite{Shi_2023}---can be seamlessly achieved using native Julia operations.

Similarly, the Organ-Specific Affine Registration pipeline converged using the custom fused GPU loss kernel. The ability of Enzyme.jl to compute exact reverse-mode gradients through the highly branched KernelAbstractions.jl code validated that complex spatial regularizations can be trained natively on the GPU without hand-writing CUDA adjoints.

\subsubsection*{Quantitative Dosimetry Refinement Performance}
The UDE-based dosimetry refinement pipeline successfully corrected approximate SPECT reconstructions into high-accuracy equivalents. To quantify performance, we benchmarked our Julia/SciML implementation against seven distinct analytical and deep learning baselines over $N=48$ $64^3$ spatial patches, using full Monte Carlo tracking as the ground truth (Table \ref{tab:dosimetry_metrics}).

\begin{table}[H]
\centering
\caption{Consolidated Dosimetry Performance Metrics ($64^3$ Patches)}
\begin{tabular}{lccc}
\toprule
\textbf{Model Architecture} & \textbf{Pearson $r$} & \textbf{MAE (norm)} & \textbf{Framework} \\ \midrule
\textbf{UDE (No-Approx)} & \textbf{0.9570} & 0.0330 & \textbf{SciML (Julia)} \\
Stabilized CNN (ResNet-32) & 0.9390 & \textbf{0.0270} & Deep Learning \\
Analytical Baseline (VSV) & 0.9120 & 0.0320 & Physics (Python) \\
PBPK-ML (Extra Trees) & 0.7723 & 0.0755 & Statistical ML \\
Spect0 (3D Residual) & 0.7471 & 1.7528 & Deep Learning \\
DblurDoseNet (3D U-Net) & 0.5566 & 0.5413 & Deep Learning \\
\bottomrule
\end{tabular}
\label{tab:dosimetry_metrics}
\end{table}

The "No-Approx UDE" implemented in Julia achieved the definitive state-of-the-art result, producing a peak Pearson correlation of \textbf{0.9570} against the Monte Carlo ground truth. This significantly outperforms both the clinical analytical standard (VSV Evaluate-Proj, $r=0.9120$) and pure deep learning frameworks (e.g., DblurDoseNet, $r=0.5566$). Figure \ref{fig:dosimetry_comparison} visually corroborates this performance, illustrating the structural accuracy of the UDE output compared to the ground truth and pure deep-learning baseline.

\begin{figure}[H]
    \centering
    \includegraphics[width=\textwidth]{elsarticle/figures_new/dosimetry_comparison_all.png}
    \caption{\textbf{Dosimetry Accuracy Comparison Across Architectures.}
    Visual assessment comparing the spatial accuracy of the "No-Approx UDE" ($r=0.9570$) against the Monte Carlo ground truth and established baselines. The physics-integrated UDE successfully models complex tissue boundary scatter, significantly reducing the unconstrained spatial approximation artifacts endemic to pure "black box" deep learning approaches.}
    \label{fig:dosimetry_comparison}
\end{figure}

Crucially, the interpretability of the UDE approach is profound. The failure of pure black-box models to exceed $r=0.60$ highlights the absolute necessity of the "Physics-in-the-loop" approach. By isolating the known mechanistic physics (primary kernel deposition, patient density map) from the complex scattering phenomena learned by the neural residual, the model captures the final 5--10\% of dose variance that traditional models either ignore (VSV) or fail to systematically learn (CNNs).

This level of precision is paramount in therapeutic nuclear medicine, where ablative doses are delivered to tumor tissue while attempting to spare healthy organs. Standard dose kernels fail to account for the heterogeneous density of tissue interfaces and "leakage" across boundaries. Furthermore, the MedImages.jl/SciML implementation maintained a 10$\times$ computational speed advantage over established Python-based analytical frameworks, offering Monte Carlo-level accuracy at speeds viable for large-scale clinical pipelines.

\subsection*{Challenge 4: Metadata Management and Clinical Fidelity}
We verified the quantitative accuracy of MedImages.jl by comparing its Standardized Uptake Value (SUV) calculations against the gold standard SimpleITK. Due to Julia's high-precision handling of temporal and clinical metadata, MedImages.jl achieves identical results for SUVbw (Body Weight) calculations, matching SimpleITK's output up to $10^{-14}$ in double precision. This ensures that the massive speedups achieved through GPU acceleration do not compromise clinical validity, a fact verifiable through the open-source benchmarks in this repository.

Furthermore, we conducted an end-to-end consistency experiment on theranostic radioligand therapy data to verify that both spatial alignment and quantitative fidelity (e.g., Standardized Uptake Values) are strictly maintained throughout complex, multi-modal spatial transformation pipelines.

Theranostic workflows integrate heterogeneous longitudinal patient imaging acquired across different modalities and reconstruction pipelines---e.g., [$^{177}$Lu]Lu-PSMA SPECT with and without attenuation correction (AC/NAC), voxel-wise absorbed dose maps, and Computed Tomography (CT) for anatomical reference. Crucially, each modality retains its own unique voxel grid, datatype, and spatial metadata (spacing, origin, orientation). To safely align or augment these images for joint neural network inference, the applied geometric transformations must preserve absolute spatial correlation between the structural anatomy and the functional activity.

We loaded these four distinct modalities into a single \texttt{BatchedMedImage} object, intrinsically preserving their individual spatial metadata. A compounded spatial sequence was then applied simultaneously as a single batched operation: (1) converting raw SPECT intensities to SUV, (2) resampling all modalities to a standardized 2.0\,mm isotropic grid, (3) applying a 45-degree rotation around the Z-axis, and (4) translating the volume by 10 voxels along the X-axis.

Quantitative analysis of the Mean SUV and total Lesion Volume inside a segmented Region of Interest (ROI) demonstrated remarkable consistency before and after these simultaneous, compound transformations. The original Mean SUV was 3.4913 with a volume of 52.55 $\text{cm}^3$, and the final transformed volume exhibited a Mean SUV of 3.4660 and a volume of 51.86 $\text{cm}^3$. The minor deviations in Mean SUV (0.73\%) and Volume (1.32\%) are fundamentally attributable to the continuous linear interpolation required for the 2.0\,mm grid resampling rather than underlying metadata deterioration.

Figure \ref{fig:transformation_consistency} visually illustrates this multi-modal batched processing and alignment stability.

\begin{figure}[H]
    \centering
    \includegraphics[width=\textwidth]{elsarticle/figures_new/transformation_consistency.png}
    \caption{\textbf{Quantitative Consistency Under Compounded Batched Spatial Transformations.}
    The composite image illustrates the simultaneous processing of a theranostic multi-modal dataset ([$^{177}$Lu]Lu-PSMA SPECT AC, SPECT NAC, Dosemap, and CT) stored within a single \texttt{BatchedMedImage}. The original volumes (left panels) underwent a compounded sequence of 2.0\,mm isotropic resampling, 45\textdegree{} Z-rotation, and X-translation to produce the augmented output (right panels).
    MedImages.jl flawlessly preserved spatial alignment between the functional distribution (SPECT/Dose) and anatomical structure (CT). Clinical metadata (SUV) remained strictly coupled, with Mean SUV and gross lesion volume varying by $<1.5\%$ (3.4913 vs 3.4660) across the entire multi-modality pipeline due to expected continuous resampling effects.}
    \label{fig:transformation_consistency}
\end{figure}


\section{Discussion}
As the field of biomedical imaging moves toward more complex 4D dynamic studies and scientific machine learning, the architectural nuances of different software frameworks become increasingly critical. Our results demonstrate that the MedImages.jl ecosystem provides a highly performant and differentiable path forward by systematically addressing the four core challenges outlined in the Introduction.

\subsection*{Addressing Challenges 1 and 2 (Volume and Speed): The Advantage of Native Fused Kernels}
The tight binding of metadata within the \texttt{MedImage} struct directly translates to the performance gains observed in our benchmarks. The 7.2$\times$ turnaround advantage observed in high-frequency training loops is primarily attributed to MedImages.jl's integration with HDF5---which bypasses the serialization costs inherent in many Python caching mechanisms---and its single-pass Fused Affine kernel. By mathematically collapsing the entire spatial transformation stack (normalization, resampling, centering) into a single CUDA execution, MedImages.jl eliminates the memory traffic and synchronization overhead associated with sequential operation queues, providing a highly efficient solution for biobank-scale deployments.

\subsection*{Addressing Challenge 4 (Metadata Management): Comparing Metadata Architectures}
Modern Python frameworks like MONAI \cite{CardosoLi2022} have introduced sophisticated mechanisms to couple metadata with image arrays. MONAI utilizes a MetaTensor structure, appending a metadata dictionary to the PyTorch Tensor. A similar design pattern appears in TorchIO \cite{perezgarcia2021torchio}, which prioritizes convenient multimodal coordination but still builds its core processing around external imaging backends. TorchIO approaches the problem through a hierarchical Subject class system. While effective for multi-modal synchronization, TorchIO is primarily a manager of third-party backends (SimpleITK, NiBabel). This dependency chain maintains a layer of indirection that can lead to memory management issues in large-scale training \cite{perezgarcia2021torchio}.

In contrast, MedImages.jl leverages Julia's type system to bind metadata directly within the \texttt{MedImage} struct. As summarized in Table \ref{tab:metadata_comparison}, this approach avoids the runtime overhead associated with dictionary lookups, contributing to the performance gains observed in our benchmarks. Furthermore, by implementing core geometric transformations directly in Julia, MedImages.jl ensures the autograd graph remains intact, a challenge in frameworks that delegate to non-differentiable compiled backends.

\begin{table}[H]
\centering
\caption{Comparative Metadata Architectures in Modern Frameworks}
\begin{tabular}{|p{2.5cm}|p{3.5cm}|p{3.5cm}|p{3.5cm}|}
\hline
\textbf{Framework} & \textbf{Metadata Storage Mechanism} & \textbf{Preservation Strategy} & \textbf{Potential Interoperability Constraints} \\
\hline
SimpleITK & Internal C++ Struct & Encapsulated in \texttt{itk::Image} & User-dependent retention (often discarded upon conversion to NumPy). \\
\hline
MONAI & MetaTensor (Dict-based) & Metadata dictionary moves with tensor & Dictionary lookup overhead; potential orientation flip during ITK interoperability. \\
\hline
TorchIO & Subject Class Attributes & Multi-modal synchronization & Potential memory buildup during extensive transformation histories. \\
\hline
MedImages.jl & Typed Julia Struct & Explicit binding within struct & N/A (Native Julia stack). \\
\hline
\end{tabular}
\label{tab:metadata_comparison}
\end{table}

\subsection*{The Primacy of Clinical and Temporal Metadata}
The lack of standardization in medical imaging pipelines can lead to a reproducibility challenge. A study on fMRI preprocessing identified over 200 unique analysis workflows across 241 studies, illustrating the impact of varied processing pipelines \cite{Esteban_2018}. In such environments, "metadata drift"---where spatial or clinical parameters become uncoupled from voxel data---can pose a clinical risk, for instance, by shifting a treatment margin in radiotherapy or invalidating quantitative interpretation in PET/CT due to lost decay tracking.
Quantitative nuclear medicine critically relies on intact clinical and temporal metadata: SUV computation depends explicitly on injection and acquisition times. In theranostic workflows, this requirement collides with the need to apply identical geometric operations across multiple co-referenced modalities (e.g., [$^{177}$Lu]Lu-PSMA SPECT AC/NAC, absorbed-dose maps, and CT).

MedImages.jl addresses this by design: the \texttt{BatchedMedImage} abstraction allows single-call, batched geometric transforms to be applied consistently across modalities without manual bookkeeping. Our cross-platform validation confirms that the SUV conversion factors extracted by MedImages.jl achieve mathematical identity with established SimpleITK benchmarks to a precision of $10^{-14}$. This numerical equivalence ensures that researchers transitioning to Julia can maintain absolute continuity with legacy clinical findings while gaining the performance required for massive biobank-scale analysis.

\subsection*{Addressing Challenge 3 (Differentiability): Navigating the Framework Ecosystem}
We do not advocate replacing established tools wholesale. Rather, we recommend choosing the framework that best matches the task at hand, especially when differentiability and complex scientific computing are core requirements.

\paragraph{When established Python frameworks remain the pragmatic choice.}
For teams with substantial legacy Python codebases, switching frameworks carries non-trivial migration costs. MONAI provides a mature library of ready-made components---pretrained segmentation networks, standardized data-loading pipelines, and extensive augmentation transforms---that can be deployed with minimal custom code. Similarly, SimpleITK offers a comprehensive catalogue of classical filters whose correctness has been validated across decades. When the research question can be answered with an existing pipeline and the primary requirement is rapid prototyping within a specific tensor framework (like PyTorch), these tools remain highly efficient.

\paragraph{When MedImages.jl offers a clear advantage for SciML.}
The Julia-based approach becomes compelling when research transcends standard convolutional network training and enters the realm of physics-informed models or custom numerical workflows \cite{RoeschGreener2023}:
\begin{itemize}
    \item \textbf{Unifying the Fractured SciML Landscape.} In Python, moving a differentiable model from PyTorch to a JAX-based differential equation solver often requires significant code rewriting or bridging overhead. Julia's multiple dispatch solves the expression problem, meaning essentially all packages are interoperable natively \cite{PalBhattacharya2024, EschleGl2023}. MedImages.jl serves as the critical entry point to make this uniquely unified SciML ecosystem accessible for medical imaging.
    \item \textbf{Novel algorithm development.} When researchers need to implement a new spatial transformation, loss function, or physics-informed model from scratch, Julia's single-language stack eliminates the need to write C++/CUDA extensions. The algorithm can be prototyped, debugged, and GPU-accelerated within the same language, with full automatic differentiation support.
    \item \textbf{Cross-domain scientific computing.} Julia's ecosystem uniquely bridges medical imaging with adjacent scientific domains---DifferentialEquations.jl for pharmacokinetic modelling, KomaMRI.jl \cite{CastilloPassiCoronado2023} for MRI pulse-sequence simulation, Flux.jl/Lux.jl \cite{pal2023lux} for deep learning---all within a single runtime. This composability enables workflows that would otherwise require fragile inter-process communication between Python, MATLAB, and C++ components.
    \item \textbf{Differentiable imaging physics and UDEs.} Tasks such as learned registration, differentiable augmentation, or physics-informed neural networks require gradients to flow through spatial transformations. The quantitative dosimetry refinement pipeline perfectly exemplifies this advantage. By formulating a Universal Differential Equation (UDE) that integrates $S_{homo}$ kernel convolutions directly within a multi-state pharmacokinetic system, our implementation achieved a $0.9570$ Pearson correlation against Monte Carlo tracking. This "Physics-in-the-loop" strategy captured the final 5--10\% of dose variance lost in pure deep learning models ($r \approx 0.60$) and surpassed established analytical methods. Because MedImages.jl's native Julia kernels integrate seamlessly with Zygote.jl and Enzyme.jl, this entire multi-modal UDE system could be trained end-to-end natively while maintaining a 10$\times$ speed advantage over equivalent Python implementations. MONAI and TorchIO, by contrast, must fall back on non-differentiable C++ backends for many geometric operations, breaking the autograd graph and making unified SciML architectures practically infeasible.
\end{itemize}

In practice, we envision a complementary workflow: established MONAI/SimpleITK pipelines for routine clinical deployment, and MedImages.jl for research frontiers where performance, differentiability, or algorithmic novelty are the primary requirements.

MedImages.jl offers a path forward---a future where imaging physics is integrated directly into the deep learning pipeline, where hardware acceleration is transparent and ubiquitous, and where the metadata that defines our physical world is preserved with the same rigor as the pixel values themselves. By bridging the divide between the physics of acquisition and the logic of inference, and by addressing the entwined challenges of \textbf{Volume} and \textbf{Speed} with a 7.2$\times$ faster turnaround than traditional caching frameworks \cite{Shi_2023}, MedImages.jl empowers researchers and clinicians to unlock the full potential of medical imaging as a precise, quantitative science. This performance is especially critical for large-scale multi-center studies and meta-analyses, where the ability to process thousands of subjects with absolute metadata integrity enables research that was previously computationally infeasible.


\subsection*{Limitations}
While MedImages.jl offers significant advantages in execution speed and differentiability, adopting the Julia ecosystem presents certain trade-offs. The most notable is the "time-to-first-plot" or compilation latency. Because Julia uses Just-In-Time compilation, the first execution of a complex spatial transformation incurs a compilation penalty, making initial startup slower than equivalent interpreted Python code. Furthermore, running full test suites or deploying models with heavy dependencies (e.g., CUDA.jl, Lux.jl, Zygote.jl) requires significant development memory, occasionally causing out-of-memory (OOM) issues on constrained hardware. Finally, while the Julia machine learning ecosystem is growing rapidly, it does not yet match the sheer volume of pre-trained models and plug-and-play architectural components available in the PyTorch/MONAI ecosystem.

\section{Conclusions}
The development of MedImages.jl addresses core computational and metadata challenges in modern medical imaging. By leveraging Julia's compilation architecture and type system, the library provides a unified framework where physical accuracy, computational performance, and end-to-end differentiability are achieved simultaneously. The native implementation of spatial transformations allows for transparent integration with SciML pipelines, offering a robust foundation for the next generation of physics-informed models in molecular imaging and quantitative theranostics.

\section*{Availability and Future Directions}
MedImages.jl is open-source and available at \url{https://github.com/JuliaHealth/MedImages.jl}. All experiments described in this paper, including performance benchmarks and differentiability validations, are reproducible via the scripts provided in the repository's \texttt{experiment} directory.

Future work will focus on:
\begin{enumerate}
    \item Integrating with advanced reconstruction frameworks (KomaMRI.jl \cite{CastilloPassiCoronado2023}, GIRFReco.jl \cite{JaffrayWu2024}) for end-to-end physics-informed learning.
    \item Implementing differentiable registration algorithms leveraging Zygote.jl.
    \item Expanding support for standardized formats like DICOM-SEG via Highdicom \cite{BridgeGorman2022} integration.
    \item Developing architectural foundations for clinical-grade AI agents (e.g., VoxelPrompt) capable of multi-modal "Medical Superintelligence" benchmarks where absolute metadata integrity is a prerequisite for safe automated diagnostics.
\end{enumerate}

\section*{Declarations}
\textbf{Ethics approval and consent to participate:} This research relies solely on retrospective analysis of fully anonymized, publicly available datasets. No new human subjects data was collected. Therefore, formal ethics committee approval and informed consent were not required. \\
\textbf{Declaration of competing interests:} The authors declare that they have no known competing financial interests or personal relationships that could have appeared to influence the work reported in this paper. \\
\textbf{Funding sources:} This research did not receive any specific grant from funding agencies in the public, commercial, or not-for-profit sectors. \\
\textbf{Declaration of generative AI and AI-assisted technologies:} During the preparation of this work, the author(s) used generative AI and AI-assisted technologies for improving style and assisting in code development. After using these tools, the authors reviewed and edited the content as needed and take full responsibility for the content of the published article.

\bibliographystyle{elsarticle-num}
\bibliography{bibl}

\end{document}
